\chapter*{Conclusion générale}
\addcontentsline{toc}{chapter}{Conclusion générale}
\markboth{\MakeUppercase{Conclusion générale}}{}

Ce Projet de Fin d'Études réalisé au sein de MEDIASSIVE m'a permis de participer à la conception, au développement, au déploiement et à l'évaluation de FinderrLink, une plateforme SaaS facilitant la mise en relation entre freelances, entreprises et agences, avec une centralisation des profils, missions, candidatures, échanges et notifications dans un environnement unique et sécurisé.

Durant cette mission, j'ai occupé le rôle de chef d'équipe stagiaire, développeur full-stack et responsable du déploiement, couvrant le développement frontend et backend, l'intégration des services, la coordination de l'équipe et la mise en ligne de la plateforme. Cette expérience m'a permis d'évoluer dans un contexte proche d'un environnement professionnel réel, avec des contraintes techniques, organisationnelles et humaines.

FinderrLink repose sur quatre rôles principaux — Freelance, Entreprise, Agence et Administrateur — chacun disposant de fonctionnalités adaptées via un système RBAC. La plateforme couvre la gestion des profils, la publication et recherche de missions, le suivi des candidatures, la messagerie, les notifications et un assistant intelligent. L'agence se distingue par son rôle hybride, lui permettant à la fois de publier des missions et de postuler à des opportunités.

Sur le plan technique, la solution s'appuie sur une architecture full-stack moderne : frontend Next.js, backend NestJS, base de données PostgreSQL, service IA en Python, et un déploiement via Docker, Nginx et VPS. Cette architecture garantit une séparation claire des responsabilités et prépare la plateforme à une évolution future.

La phase d'évaluation, menée à travers des tests manuels, des validations backend et un test de charge k6, confirme le bon fonctionnement général de FinderrLink tout en identifiant des pistes d'amélioration. Ce stage a renforcé mes compétences en développement web, déploiement et gestion de projet. FinderrLink reste une plateforme évolutive, avec des perspectives incluant les tests automatisés, la visioconférence, le support multilingue, l'IA agentique et une application mobile native.

